% GENERATED FILE. Edit canonical lessons, then run npm run generate:books.
% canonical-source-hash: fnv1a64:07c9f4fe330c3c7b
% canonical-lessons: BN-C420-bikri, BN-C420-pherioyala, BN-C420-dalal, BN-C420-hat, BN-C420-bikono

\chapter{Sales, Hawkers and Markets}
\label{ch:bn-sales-hawkers-and-markets}
\hlchaptermodality{420}

\begin{chapteropening}
\textbf{By the end of this chapter:} \emph{I can talk about selling at a village market, to a hawker or through a broker.}

Complete the last lesson of chapter 420: I can talk about selling at a village market, to a hawker or through a broker.
\end{chapteropening}

\section[\texorpdfstring{\bn{বিক্রি} (bikri)}{বিক্রি (bikri)}]{\bn{বিক্রি} (bikri) — a sale; selling}

\label{lesson:BN-C420-bikri}



\begin{quote}
Before the new one: say the Bengali for a pledge (for a loan), then the Bengali for one's means.
\end{quote}

\subsection*{You'll want to know: \bn{বিক্রি}}
\textbf{\bn{বিক্রি}} — \emph{bikri} — \textquotedblleft{}a sale; selling\textquotedblright{}.

\begin{quote}
\bn{শীতে} \bn{কম্বলের} \bn{বিক্রি} \bn{ভালো} \bn{হয়।}
\end{quote}

\begin{quote}
\emph{shīte kômbôler bikri bhālo hôy.}
\end{quote}

\begin{quote}
Blankets sell well in winter.
\end{quote}

\bn{বিক্রি} \bn{করা}, to sell; \bn{বিক্রি} \bn{হওয়া}, to be sold.

\begin{cousinweb}
\textbf{\bn{বিক্রি}} (\emph{bikri}) goes back to Sanskrit \emph{vikrayá}, \textquotedblleft{}a sale\textquotedblright{} (CDIAL 11638), under which CDIAL lists the older Bengali \emph{bikā} and \emph{biki}; \bn{বিক্রি} is the half-learned form in use today.
\end{cousinweb}

\begin{grammarlens}[title={What you've built}]
One more word for getting around.
\end{grammarlens}

\subsection*{Your turn}
\begin{itemize}
\raggedright
  \item \emph{Say it:} \emph{bikri}
  \item \emph{Say it:} \emph{bikri}, once more
  \item \emph{Say it:} \emph{shīte kômbôler bikri bhālo hôy.} — Blankets sell well in winter.
  \item \emph{Recall:} say the Bengali for a pledge (for a loan), then the Bengali for one's means, then say \emph{bikri} again
\end{itemize}

\begin{tcolorbox}[breakable,colback=teal!4,colframe=teal!35!black,title={Before you move on}]
What does \bn{বিক্রি} mean? (\textquotedblleft{}A sale; selling\textquotedblright{}.) Say it once more.
\end{tcolorbox}
\section[\texorpdfstring{\bn{ফেরিওয়ালা} (pherioyālā)}{ফেরিওয়ালা (pherioyālā)}]{\bn{ফেরিওয়ালা} (pherioyālā) — a hawker; a street seller}

\label{lesson:BN-C420-pherioyala}



\begin{quote}
Before the new one: say the Bengali for one's means, then the Bengali for a sale (selling).
\end{quote}

\subsection*{You'll want to know: \bn{ফেরিওয়ালা}}
\textbf{\bn{ফেরিওয়ালা}} — \emph{pherioyālā} — \textquotedblleft{}a hawker; a street seller\textquotedblright{}.

\begin{quote}
\bn{দুপুরে} \bn{গলিতে} \bn{ফেরিওয়ালা} \bn{হাঁকছিল।}
\end{quote}

\begin{quote}
\emph{dupure gôlite pherioyālā hā̃kchhilô.}
\end{quote}

\begin{quote}
At midday a hawker was crying his wares in the lane.
\end{quote}

\begin{cousinweb}
\textbf{\bn{ফেরিওয়ালা}} (\emph{pherioyālā}) joins \bn{ফেরি}, \textquotedblleft{}going round selling\textquotedblright{}, from the root of \bn{ফেরা}, to turn (CDIAL 9078), to -\bn{ওয়ালা}, the ending for one who deals in a thing: the man who goes round.
\end{cousinweb}

\begin{grammarlens}[title={What you've built}]
One more word for getting around.
\end{grammarlens}

\subsection*{Your turn}
\begin{itemize}
\raggedright
  \item \emph{Say it:} \emph{pherioyālā}
  \item \emph{Say it:} \emph{pherioyālā}, once more
  \item \emph{Say it:} \emph{dupure gôlite pherioyālā hā̃kchhilô.} — At midday a hawker was crying his wares in the lane.
  \item \emph{Recall:} say the Bengali for one's means, then the Bengali for a sale (selling), then say \emph{pherioyālā} again
\end{itemize}

\begin{tcolorbox}[breakable,colback=teal!4,colframe=teal!35!black,title={Before you move on}]
What does \bn{ফেরিওয়ালা} mean? (\textquotedblleft{}A hawker; a street seller\textquotedblright{}.) Say it once more.
\end{tcolorbox}
\section[\texorpdfstring{\bn{দালাল} (dālāl)}{দালাল (dālāl)}]{\bn{দালাল} (dālāl) — a broker; a middleman}

\label{lesson:BN-C420-dalal}



\begin{quote}
Before the new one: say the Bengali for a sale (selling), then the Bengali for a hawker.
\end{quote}

\subsection*{You'll want to know: \bn{দালাল}}
\textbf{\bn{দালাল}} — \emph{dālāl} — \textquotedblleft{}a broker; a middleman\textquotedblright{}.

\begin{quote}
\bn{দালাল} \bn{বাসার} \bn{জন্য} \bn{অগ্রিম} \bn{চাইল।}
\end{quote}

\begin{quote}
\emph{dālāl bāsār jônyô ôgrim chāilô.}
\end{quote}

\begin{quote}
The broker asked for an advance for the flat.
\end{quote}

\begin{cousinweb}
\textbf{\bn{দালাল}} (\emph{dālāl}) is usually traced to Arabic \emph{dallāl}, \textquotedblleft{}a broker, a guide\textquotedblright{}, through Persian; the word often carries a hint of distrust.
\end{cousinweb}

\begin{grammarlens}[title={What you've built}]
One more word for getting around.
\end{grammarlens}

\subsection*{Your turn}
\begin{itemize}
\raggedright
  \item \emph{Say it:} \emph{dālāl}
  \item \emph{Say it:} \emph{dālāl}, once more
  \item \emph{Say it:} \emph{dālāl bāsār jônyô ôgrim chāilô.} — The broker asked for an advance for the flat.
  \item \emph{Recall:} say the Bengali for a sale (selling), then the Bengali for a hawker, then say \emph{dālāl} again
\end{itemize}

\begin{tcolorbox}[breakable,colback=teal!4,colframe=teal!35!black,title={Before you move on}]
What does \bn{দালাল} mean? (\textquotedblleft{}A broker; a middleman\textquotedblright{}.) Say it once more.
\end{tcolorbox}
\section[\texorpdfstring{\bn{হাট} (hāṭ)}{হাট (hāṭ)}]{\bn{হাট} (hāṭ) — a village market (on set days)}

\label{lesson:BN-C420-hat}



\begin{quote}
Before the new one: say the Bengali for a hawker, then the Bengali for a broker.
\end{quote}

\subsection*{You'll want to know: \bn{হাট}}
\textbf{\bn{হাট}} — \emph{hāṭ} — \textquotedblleft{}a village market (on set days)\textquotedblright{}.

\begin{quote}
\bn{আমাদের} \bn{গ্রামে} \bn{বুধবার} \bn{হাট} \bn{বসে।}
\end{quote}

\begin{quote}
\emph{āmāder grāme budhbār hāṭ bôse.}
\end{quote}

\begin{quote}
In our village the market is held every Wednesday.
\end{quote}

A \bn{বাজার} opens every day; a \bn{হাট} \bn{বসে}, sits, on its own days of the week.

\begin{cousinweb}
\textbf{\bn{হাট}} (\emph{hāṭ}) is inherited: CDIAL lists Bengali \emph{hāṭ}, \textquotedblleft{}a shop, a market\textquotedblright{}, under Sanskrit \emph{haṭṭa}, \textquotedblleft{}a market, a fair\textquotedblright{} (CDIAL 13944).
\end{cousinweb}

\begin{grammarlens}[title={What you've built}]
One more word for getting around.
\end{grammarlens}

\subsection*{Your turn}
\begin{itemize}
\raggedright
  \item \emph{Say it:} \emph{hāṭ}
  \item \emph{Say it:} \emph{hāṭ}, once more
  \item \emph{Say it:} \emph{āmāder grāme budhbār hāṭ bôse.} — In our village the market is held every Wednesday.
  \item \emph{Recall:} say the Bengali for a hawker, then the Bengali for a broker, then say \emph{hāṭ} again
\end{itemize}

\begin{tcolorbox}[breakable,colback=teal!4,colframe=teal!35!black,title={Before you move on}]
What does \bn{হাট} mean? (\textquotedblleft{}A village market (on set days)\textquotedblright{}.) Say it once more.
\end{tcolorbox}
\section[\texorpdfstring{\bn{বিকোনো} (bikono)}{বিকোনো (bikono)}]{\bn{বিকোনো} (bikono) — to be sold; to sell (for a price)}

\label{lesson:BN-C420-bikono}



\begin{quote}
Before the new one: say the Bengali for a broker, then the Bengali for a village market.
\end{quote}

\subsection*{You'll want to know: \bn{বিকোনো}}
\textbf{\bn{বিকোনো}} — \emph{bikono} — \textquotedblleft{}to be sold; to sell (for a price)\textquotedblright{}.

\begin{quote}
\bn{হাটে} \bn{মাছটা} \bn{ভালো} \bn{দামে} \bn{বিকোল।}
\end{quote}

\begin{quote}
\emph{hāṭe māchhṭā bhālo dāme bikolô.}
\end{quote}

\begin{quote}
At the market the fish sold for a good price.
\end{quote}

Past (he/she …): \bn{সেই} \bn{মাছ} \bn{ভালো} \bn{দামে} \bn{বিকোল}, \emph{sei māchh bhālo dāme bikolô}.

The goods are the subject: \bn{আম} \bn{বিকোয়}, mangoes sell.

\begin{cousinweb}
\textbf{\bn{বিকোনো}} (\emph{bikono}) is inherited: CDIAL lists Bengali \emph{bikā}, \textquotedblleft{}to sell\textquotedblright{}, under Sanskrit \emph{vikrīyatē}, \textquotedblleft{}is sold\textquotedblright{} (CDIAL 11642).
\end{cousinweb}

\begin{grammarlens}[title={What you've built}]
One more word for getting around.
\end{grammarlens}

\subsection*{Your turn}
\begin{itemize}
\raggedright
  \item \emph{Say it:} \emph{bikono}
  \item \emph{Say it:} \emph{bikono}, once more
  \item \emph{Say it:} \emph{hāṭe māchhṭā bhālo dāme bikolô.} — At the market the fish sold for a good price.
  \item \emph{Recall:} say the Bengali for a broker, then the Bengali for a village market, then say \emph{bikono} again
\end{itemize}

\begin{tcolorbox}[breakable,colback=teal!4,colframe=teal!35!black,title={Before you move on}]
What does \bn{বিকোনো} mean? (\textquotedblleft{}To be sold; to sell (for a price)\textquotedblright{}.) Say it once more.
\end{tcolorbox}
